\section{Estado de la implementación}
\label{sec:resultados}

La Tabla~\ref{tab:estado} resume el estado de cada funcionalidad a la
fecha de esta entrega, distinguiendo entre funcionalidades implementadas,
parcialmente implementadas y pendientes.

\begin{table}[t]
  \caption{Estado de las funcionalidades en la primera entrega}
  \label{tab:estado}
  \centering
  \begin{tabular}{p{0.58\columnwidth}p{0.28\columnwidth}}
    \toprule
    Funcionalidad & Estado \\
    \midrule
    Registrar solicitud (RF1) & Implementada \\
    Atender por prioridad (RF2) & Implementada \\
    Buscar por identificador (RF3) & Implementada \\
    Cancelar solicitud (RF4) & Implementada \\
    Deshacer última atención (RF5) & Pendiente \\
    Consultar historial (RF6) & Implementada \\
    Interfaz de consola & Implementada \\
    Validación de entradas & Parcial \\
    Pruebas del sistema & Parcial \\
    Índice por identificador en $O(1)$ & Pendiente \\
    Análisis experimental de desempeño & Pendiente \\
    \bottomrule
  \end{tabular}
\end{table}

\textbf{Implementadas.} Cinco de los seis requerimientos funcionales están
disponibles en la clase \texttt{TriageSystem}: \texttt{registerRequest},
\texttt{serveNext}, \texttt{findById}, \texttt{cancelRequest} y
\texttt{getHistory}. La interfaz de consola (\texttt{Main}) expone todas
estas operaciones mediante un menú. Las estructuras de soporte
(\texttt{LinkedQueue}, \texttt{LevelPriorityQueue} y \texttt{LinkedStack})
son implementación propia; de la biblioteca estándar de Java solo se usa
\texttt{ArrayList} para devolver el historial como lista.

\textbf{Parcialmente implementadas.} La validación de entradas rechaza
identificadores vacíos o repetidos entre las solicitudes pendientes,
urgencias fuera del rango 1--5 y urgencias no numéricas, pero aún no valida
el nombre de la solicitud. Las pruebas se realizan mediante un programa de
prueba manual (\texttt{BasicTests.java}) que imprime el resultado esperado
junto al obtenido, sin aserciones automáticas. El escenario principal
registra cuatro solicitudes, verifica el orden de atención por urgencia y
el desempate por orden de llegada entre dos solicitudes de urgencia 1,
cancela una solicitud pendiente y comprueba que deja de encontrarse y de
contarse, y confirma el orden del historial. Los casos borde cubren una
urgencia fuera de rango, un identificador repetido, la cancelación de un
identificador inexistente, la reutilización del identificador de una
solicitud ya atendida y la atención con la cola vacía.

\textbf{Pendientes.} Para la siguiente entrega quedan deshacer la última
atención (RF5), cuyo diseño se apoya en la pila ya implementada; un índice
por identificador que reduzca la búsqueda y la cancelación a tiempo
esperado constante; y la migración de las pruebas manuales a pruebas
automatizadas. Para la entrega final queda un análisis experimental de
tiempos de ejecución sobre volúmenes de datos mayores.

No se reportan resultados experimentales de desempeño en esta entrega,
dado que el énfasis de la primera versión es la corrección funcional y la
justificación del diseño; el análisis de complejidad teórica se presentó
en la Sección~\ref{sec:metodologia}.
